% Written by paper/cpc/scripts/corpus31_results.py from data/corpus31-results.json and data/corpus31/manifest.toml.
\begin{longtable}{@{}p{0.20\linewidth}p{0.27\linewidth}p{0.14\linewidth}p{0.29\linewidth}@{}}
\caption{Entries of the CON 3.1 corpus whose outcome no clause decides by name, each with the reading that its manifest entry records.}\label{tab:corpus31-open}\\
\toprule
File & Outcome; basis & Clauses & Reading\\
\midrule
\endfirsthead
\toprule
File & Outcome; basis & Clauses & Reading\\
\midrule
\endhead
\bottomrule
\endfoot
\path{v2_minor1_velocity_no_time.convel} & accept, 1 frame; open & CON-units.8, CON-units.11 & Major 2 with con\_spec\_minor 1. CON-units.8 reads version 2 velocities in the natural time unit, and CON-units.11 rejects a frame of minor 1 or later without units.time. The reader takes a minor to refine major 3 only, so it reads the velocities in the natural time unit.\\
\path{v31_units_compound.convel} & accept, 1 frame; open & CON-units.1, CON-units.2, CON-units.16 & Products, quotients and parentheses of the matching dimensions (Units). The energy is written kJ / mol; the text does not say whether white space may stand inside an expression, and the reader removes it.\\
\path{v32_unknown_section_arity.con} & accept, 1 frame; optional & CON-line2.7, CON-section.8, CON-section.14 & Frame minor 2 is newer than the reader's minor 1 and section\_arity gives stress 6 value fields, so the reader may skip the stress section, and it does.\\
\path{v33_unknown_section_arity.con} & accept, 1 frame; optional & CON-line2.7, CON-section.8, CON-section.14 & Frame minor 3 is newer than the reader's minor 1 and section\_arity gives stress 6 value fields, so the reader may skip the stress section, and it does.\\
\path{v31_two_blanks_before_section.convel} & \texttt{IncompleteSection}, line 16; class-open & CON-section.2, CON-section.4, CON-section.15 & Two blank lines before the velocities section: the first opens it, the second is its symbol line, and the line Pt fails as its label. Outside validate true no clause names the class of a wrong label; the reader reports the declared section as missing.\\
\path{v31_unicode_blank_before_section.convel} & accept, 1 frame; open & CON-section.2, CON-file.4 & The line before the velocities section holds U+2003 only. The reader takes a line of Unicode White\_Space characters to be blank before a section or a legacy velocity block. The ignored end of the file is one such line, then lines of spaces, tabs, form feeds and carriage returns only. The text does not define whitespace or blank otherwise. So the line is blank and opens the declared section.\\
\path{v31_spreads_negative.con} & \texttt{ValidationError}, line 17; class-open & CON-section.20 & Each spreads value must be >= 0 in every mode; the clause names no class, and the reader reports ValidationError.\\
\path{v31_section_label_mismatch.convel} & \texttt{IncompleteSection}, line 16; class-open & CON-section.15 & The section label reads Velocity of Component 1 without validate. CON-section.15 applies in every mode and names no class; the reader reports the declared section as missing.\\
\path{v31_label_velocities_no_index.convel} & \texttt{IncompleteSection}, line 16; class-open & CON-section.15 & The label of component 1 of the velocities section reads Velocities of Component; the reader reports the declared section as missing.\\
\path{v31_label_velocities_xyz.convel} & \texttt{IncompleteSection}, line 16; class-open & CON-section.15 & The label of component 1 of the velocities section reads Velocities of Component xyz; the reader reports the declared section as missing.\\
\path{v31_label_velocities_1_xyz.convel} & \texttt{IncompleteSection}, line 16; class-open & CON-section.15 & The label of component 1 of the velocities section reads Velocities of Component 1 xyz; the reader reports the declared section as missing.\\
\path{v31_label_velocities_0.convel} & \texttt{IncompleteSection}, line 16; class-open & CON-section.15 & The label of component 1 of the velocities section reads Velocities of Component 0; the reader reports the declared section as missing.\\
\path{v31_label_velocities_2.convel} & \texttt{IncompleteSection}, line 16; class-open & CON-section.15 & The label of component 1 of the velocities section reads Velocities of Component 2; the reader reports the declared section as missing.\\
\path{v31_label_velocities_01.convel} & \texttt{IncompleteSection}, line 16; class-open & CON-section.15 & The label of component 1 of the velocities section reads Velocities of Component 01; the reader reports the declared section as missing.\\
\path{v31_coord_label_mismatch.con} & \texttt{ValidationError}, line 11; open & CON-validate.4 & The coordinate label reads Coords of Component 1 without validate. CON-validate.4 lists the coordinate label check under validate true, and no clause asks for it in other modes; the reader checks it in every mode.\\
\path{v31_label_coordinates_2.con} & \texttt{ValidationError}, line 11; class-open & CON-count.8 & The coordinate label of component 1 reads Coordinates of Component 2, and component numbers start at 1.\\
\path{v31_label_coordinates_1_xyz.con} & \texttt{ValidationError}, line 11; open & CON-validate.4 & The coordinate label of component 1 reads Coordinates of Component 1 xyz without validate; as for \path{v31_coord_label_mismatch.con}, the reader checks the label in every mode.\\
\path{v31_sections_out_of_order.convel} & \texttt{IncompleteSection}, line 16; open & CON-section.1, CON-section.15 & sections lists forces then velocities and the blocks are written velocities then forces. No clause fixes the order of the section blocks; the reader reads them in sections order, so the first label fails.\\
\path{v31_blank_symbol_velocities.convel} & accept, 1 frame; open & CON-section.4, CON-validate.4 & The symbol line of the velocities block is blank. CON-validate.4 checks section symbols under validate true, and no clause says what a symbol line holds in other modes; the reader does not check it.\\
\path{v31_section_row_id_mismatch.convel} & accept, 1 frame; optional & CON-section.17 & A velocity row's id differs from its coordinate row and validate is absent; the reader may ignore the mismatch, and it does.\\
\path{f26_cell_10_10_170.con} & \texttt{ValidationError}, line 4; class-open & CON-cell.6 & Round-2 report B, item F26: validate true, lengths 10 10 10 and angles 10 10 170; the radicand of the c vector of CON-cell.6 is negative, so no real cell exists; no clause names the class, and the reader rejects the angles line as ValidationError in every mode.\\
\path{gram_170_v1.con} & \texttt{ValidationError}, line 4; class-open & CON-cell.6 & Angles 170 170 170 in a version 1 frame: the radicand of the c vector of CON-cell.6 is negative, so no real cell exists; no clause names the class, and the reader rejects the angles line as ValidationError in every mode.\\
\path{gram_170_v2.con} & \texttt{ValidationError}, line 4; class-open & CON-cell.6 & Angles 170 170 170 in a version 2 frame: the radicand of the c vector of CON-cell.6 is negative, so no real cell exists; no clause names the class, and the reader rejects the angles line as ValidationError in every mode.\\
\path{gram_170_v30.con} & \texttt{ValidationError}, line 4; class-open & CON-cell.6 & Angles 170 170 170 in a version 3.0 frame: the radicand of the c vector of CON-cell.6 is negative, so no real cell exists; no clause names the class, and the reader rejects the angles line as ValidationError in every mode.\\
\path{gram_170_v31_validate.con} & \texttt{ValidationError}, line 4; class-open & CON-cell.6 & Angles 170 170 170 in a minor 1, validate true frame: the radicand of the c vector of CON-cell.6 is negative, so no real cell exists; no clause names the class, and the reader rejects the angles line as ValidationError in every mode.\\
\path{gram_60_60_120_1.con} & \texttt{ValidationError}, line 4; class-open & CON-cell.6 & Angles 60 60 120.1, validate absent: the radicand of the c vector of CON-cell.6 is negative, so no real cell exists; no clause names the class, and the reader rejects the angles line as ValidationError in every mode.\\
\path{gram_60_60_120_1_validate.con} & \texttt{ValidationError}, line 4; class-open & CON-cell.6 & Angles 60 60 120.1, validate true: the radicand of the c vector of CON-cell.6 is negative, so no real cell exists; no clause names the class, and the reader rejects the angles line as ValidationError in every mode.\\
\path{v31_version_3_5.con} & \texttt{InvalidMetadataJson}, line 2; open & CON-line2.4, CON-line2.6 & con\_spec\_version 3.5 is a JSON number that names no integer major. CON-line2.6 makes an unknown major UnsupportedVersion, and the error table makes a reserved key of the wrong type InvalidMetadataJson; the reader takes a major to be an integer and reports InvalidMetadataJson.\\
\path{v2_json_depth_127.con} & accept, 1 frame; open & CON-line2.1 & Line 2 nests 127 deep counting the outer object. No clause bounds the nesting depth of line 2, and RFC 8259 section 9 lets a parser set one. The reader accepts 127 levels counting the outer object and rejects deeper nesting.\\
\path{v2_json_depth_128.con} & \texttt{InvalidMetadataJson}, line 2; open & CON-line2.1, CON-line2.2 & Line 2 nests 128 deep counting the outer object. No clause bounds the nesting depth of line 2, and RFC 8259 section 9 lets a parser set one. The reader accepts 127 levels counting the outer object and rejects deeper nesting.\\
\path{d06_lattice_1e999.con} & \texttt{InvalidMetadataJson}, line 2; open & CON-cell.5, CON-line2.13 & Probe D6: 1e999 is a JSON number whose binary64 value is infinite. CON-cell.5 asks for finite numbers, and the error table makes an Inf value NonFiniteValue. A reserved key must hold the type this document gives it (CON-line2.13), and the reader reports a reserved key with a number that is not a finite binary64 as InvalidMetadataJson on line 2.\\
\path{d10_json_depth_129.con} & \texttt{InvalidMetadataJson}, line 2; open & CON-line2.1, CON-line2.2 & Probe D10: line 2 nests 129 deep counting the outer object. No clause bounds the nesting depth of line 2, and RFC 8259 section 9 lets a parser set one. The reader accepts 127 levels counting the outer object and rejects deeper nesting.\\
\path{d11_validate_pbc_object.con} & \texttt{InvalidMetadataJson}, line 2; open & CON-cell.4, CON-validate.4 & Probe D11: validate true and pbc is an object. The error table makes a reserved key of the wrong type InvalidMetadataJson in every mode, and CON-validate.4 names ValidationError for a bad reserved-key type under validate true. The reader applies the check of every mode first.\\
\path{d12_length_trailing_space.con} & accept, 1 frame; open & CON-units.1, CON-units.16 & Probe D12: units.length is angstrom with a trailing space. The text does not say whether white space may stand inside a unit expression; the reader removes it.\\
\path{d13_blank_symbol_energies.con} & accept, 1 frame; open & CON-section.4, CON-validate.4 & Probe D13: the symbol line of the energies block is blank. As for \path{v31_blank_symbol_velocities.convel}, the reader does not check section symbols outside validate true.\\
\path{d15_json_depth_100001.con} & \texttt{InvalidMetadataJson}, line 2; open & CON-line2.1, CON-line2.2 & Probe D15: line 2 nests 100001 deep, and the reader rejects it without exhausting its stack. No clause bounds the nesting depth of line 2, and RFC 8259 section 9 lets a parser set one. The reader accepts 127 levels counting the outer object and rejects deeper nesting.\\
\path{v31_type_count_2pow64.con} & \texttt{InvalidCount}, line 7; open & CON-error.3 & The type count is 2\^{}64. No clause bounds a count; the reader reads counts as integers in 0 through 2\^{}64-1 and reports a larger one as InvalidCount.\\
\path{v31_gzip_by_content.con} & accept, 1 frame; optional & CON-compress.3 & Gzip bytes under a plain .con name; readers should detect gzip by its magic bytes, and this one does.\\
\path{v31_bom.con} & \texttt{ValidationError}, line 1; class-open & CON-file.2 & Begins with the UTF-8 byte-order mark EF BB BF; the clause names no class.\\
\path{v31_utf8_line1.con} & \texttt{ValidationError}, line 1; class-open & CON-file.1 & Line 1 holds the bytes C0 AF, which are not UTF-8; the clause names no class.\\
\path{v31_utf8_number_token.con} & \texttt{InvalidNumberFormat}, line 12; class-open & CON-file.1, CON-num.8 & The byte C0 inside the first coordinate of line 12 is not UTF-8, and the token is not a decimal. CON-num.8 names InvalidNumberFormat and CON-file.1 names no class; the reader reports InvalidNumberFormat.\\
\path{v31_utf8_separator_line.convel} & \texttt{ValidationError}, line 14; class-open & CON-file.1, CON-section.2 & Line 14, where the velocities section needs its blank line, holds the bytes C0 AF, which are not UTF-8; neither clause names a class.\\
\path{v31_unicode_tail.con} & \texttt{IncompleteHeader}, line 16; open & CON-frame.1, CON-frame.3 & Lines of U+00A0 U+3000 and of U+2003 follow the last frame. The reader takes a line of Unicode White\_Space characters to be blank before a section or a legacy velocity block. The ignored end of the file is one such line, then lines of spaces, tabs, form feeds and carriage returns only. The text does not define whitespace or blank otherwise. The second line holds U+2003, so the two lines are lines 1 and 2 of a further frame, which ends before its line 3.\\
\path{ff_tail_610_frames.con} & accept, 610 frames; open & CON-frame.1 & 610 legal frames and a last line holding a form feed. The reader takes a line of Unicode White\_Space characters to be blank before a section or a legacy velocity block. The ignored end of the file is one such line, then lines of spaces, tabs, form feeds and carriage returns only. The text does not define whitespace or blank otherwise. So the form feed line is the ignored end of the file.\\
\path{whitespace_only.con} & \texttt{InvalidNumberFormat}, line 3; open & CON-frame.1, CON-num.8 & Lines 1 and 2 hold a space and a tab, line 3 holds U+3000. The reader takes a line of Unicode White\_Space characters to be blank before a section or a legacy velocity block. The ignored end of the file is one such line, then lines of spaces, tabs, form feeds and carriage returns only. The text does not define whitespace or blank otherwise. Line 3 holds U+3000, so the file holds a frame whose line 3 is not a decimal.\\
\path{v31_zero_atom_unicode_blank.con} & \texttt{InvalidCount}, line 21; open & CON-count.1 & A zero-atom frame whose lines 8 and 9 hold U+3000 and U+00A0. The reader takes white space inside a line to be spaces and tabs, the separators of CON-token.3 and the trailing characters of CON-file.4; the text does not define whitespace or blank otherwise. So line 8 holds a token, which CON-count.1 makes InvalidCount.\\
\path{v31_label_case_coordinates.con} & accept, 1 frame; open & CON-file.5, CON-validate.4 & The coordinate label has leading spaces, a tab between words and upper case. CON-file.5 collapses runs of spaces; the reader compares labels word by word, splitting at spaces and tabs, with ASCII letters folded.\\
\path{empty.con} & accept, 0 frames; open & CON-frame.1 & Zero bytes. No clause speaks to a file without a frame; the reader returns zero frames.\\
\end{longtable}
